\section{Alcance, delimitaciones y estructura del documento}

\subsection{Alcance}

Esta tesis desarrolla y valida un sistema híbrido a escala de laboratorio (mesocosmos) para la estimación dinámica de emisiones de \ce{CO2} en la bioconversión con \textit{Hermetia illucens}. El alcance técnico comprende: (i) la instrumentación de un nodo IoT con sensores NDIR de bajo costo para \ce{CO2} y \ce{CH4}; (ii) la formulación de un modelo mecanicista basado en EDO siguiendo el marco Dynamic Energy Budget (DEB); (iii) la integración de un calibrador neuronal parámetrico (MLP) que refine las predicciones del modelo DEB; y (iv) la validación del sistema mediante métricas de regresión, detección de anaerobiosis e inventario dinámico de carbono.

El sistema no pretende sustituir el juicio técnico del operador, sino actuar como herramienta de soporte a la decisión. Su propósito es proporcionar a productores información cuantitativa para correlacionar decisiones operativas (densidad, sustrato, aireación) con la huella de carbono resultante.

\subsection{Delimitaciones}

\begin{itemize}
    \item \textbf{Variables de GEI:} No se mide directamente \ce{N2O} en esta etapa por limitaciones de costo en sensores de bajo costo, aunque la arquitectura modular del sistema está diseñada para permitir su integración futura, reconociendo su importancia en el balance global de GEI.
    
    \item \textbf{Horizonte temporal experimental:} La validación experimental se restringe a la fase de crecimiento larval temprano (días 0--7 del ciclo) con dos dietas: sustrato de pollo estándar y mezcla de residuos agroindustriales (pulpa de café, cáscaras de plátano y naranja). Los estadios tardíos (L5/prepupa) y la fase de metamorfosis quedan fuera del alcance experimental.
    
    \item \textbf{Condiciones ambientales:} El protocolo experimental opera en régimen isotermo ($30 \pm 2^\circ\text{C}$) con humedad relativa controlada (60--80\%). La modulación térmica de Sharpe--Schoolfield se valida en este rango; la variación térmica intencional queda como trabajo futuro.
    
    \item \textbf{Escala espacial:} Los reactores batch tienen volumen útil de 0.7--0.9 L con 700 larvas y 250 g de sustrato. Los resultados son extrapolables a escala piloto con reservas metodológicas, pero no a escala industrial sin validación adicional.
    
    \item \textbf{Arquitectura de ML:} El sistema explora la estimación directa de biomasa mediante RNA (objetivo específico 3, dimensión exploratoria) y la arquitectura híbrida DEB--MLP calibrador parámetrico. La RNA directa se implementa y valida como baseline; el MLP calibrador es la arquitectura aprobada para la operación del sistema.
\end{itemize}

\section{Ajustes metodológicos respecto al anteproyecto aprobado}

El objetivo específico 3 del proyecto doctoral aprobado planteó originalmente: 
\textit{``Desarrollar un algoritmo de estimación de estados mediante redes neuronales 
que infiera la biomasa larval a partir de datos operativos, integrando esta variable 
como entrada dinámica al modelo mecanicista (DEB) para mejorar la precisión de la 
estimación de CO$_2$''}. 

Durante la fase de implementación, esta vía fue explorada rigurosamente mediante 
una arquitectura de red neuronal directa (RNA-B) entrenada con mediciones destructivas 
de biomasa como supervisión. Los resultados de validación cruzada por réplica 
(LORO, $k=6$ folds) revelaron un fallo sistemático: la RNA-B alcanzó un MAPE de 
entrenamiento del $7{,}8$\,\%pero divergió a $>40\,$\% en validación, emitiendo 
predicciones no físicas (crecimiento negativo en ventanas donde el metabolismo 
larval obliga $dB/dt \geq 0$). 

Este resultado negativo no constituye una incapacidad del algoritmo, sino una 
propiedad estructural del problema inverso. El operador forward del modelo DEB, 
$\mathcal{F}: (S,E,B,L) \mapsto \dot{m}_{\mathrm{CO_2}}$, no es inyectivo en ventanas 
cortas de observación ($\tau \leq 5$ min): múltiples trayectorias del espacio de 
estados biológicos pueden producir idénticas tasas respiratorias instantáneas. 
Con un corpus de entrenamiento destructivo limitado ($N \approx 30$ observaciones 
por campaña), la red neuronal carece de información suficiente para desambiguar 
estas trayectorias metabólicamente equivalentes, incurriendo en sobreajuste 
estructural a las réplicas de entrenamiento y pérdida de generalización inter-réplica.

Frente a esta evidencia, se adoptó una decisión metodológica deliberada, discutida 
y validada por la dirección de tesis: preservar el componente de redes neuronales 
exigido por el objetivo específico 3, pero reconfigurar su función epistemológica. 
En lugar de intentar invertir $\mathcal{F}^{-1}$ directamente desde el espacio de 
observables —vía estadísticamente mal condicionada—, se implementó un 
\textbf{calibrador neuronal paramétrico} (MLP) que opera como corrector de escala 
sobre el subconjunto de parámetros de primera sensibilidad del DEB: 
$\{\tau_h, \rho_{\mathrm{conv}}, T_{\mathrm{opt}}\}$. El MLP nunca recibe biomasa como 
supervisión; su señal de error es el residuo de CO$_2$ $\varepsilon_k = s_k - J_k^{\mathrm{DEB}}$, 
delegando la dinámica biológica al modelo mecanicista forward y reduciendo el espacio 
de corrección a una dimensión manejable ($\mathbb{R}^3$) donde la identificabilidad 
biofísica está garantizada.

En consecuencia, el objetivo específico 3 se cumple en su \textbf{dimensión 
exploratoria}: se probó, documentó y descartó formalmente la vía de estimación 
directa de biomasa mediante RNA pura, aportando evidencia empírica sobre los 
límites de las redes neuronales de caja negra en sistemas biológicos con muestreo 
destructivo limitado. Simultáneamente, se cumple en su \textbf{dimensión 
operativa}: la red neuronal integrada (MLP corrector) mejora la precisión del 
sistema híbrido hasta MAPE $= 8,3\,$\% y $R^2 = 0,78$, superando ampliamente la 
línea base estática y preservando la interpretabilidad fisiológica de cada 
parámetro. Por tanto, se considera cumplido el objetivo específico 3 en su 
totalidad: se desarrolló un algoritmo de estimación de estados mediante redes 
neuronales (el MLP calibrador), que infiere indirectamente el estado del sistema 
a través de la corrección paramétrica y mejora la precisión del modelo DEB para 
la estimación de CO$_2$. La biomasa estructural $B(t)$ se recupera indirectamente 
como salida del DEB re-simulado con parámetros corregidos, no como predicción 
directa de la red.

La Sección 4.6 documenta la evaluación comparativa formal entre la RNA-B 
(explorada y descartada) y el MLP calibrador (adoptado), garantizando la 
trazabilidad completa del ajuste metodológico.

\section{Estructura del documento}

El documento se organiza en seis capítulos principales más anexos técnicos:

\begin{description}
    \item[Capítulo 1] \textbf{Introducción General.} Situación del problema, contexto de residuos agroindustriales y cambio climático, bioconversión con BSF, planteamiento del problema, pregunta de investigación, hipótesis, objetivos y alcance.
    
    \item[Capítulo 2] \textbf{Instrumentación IoT y Caracterización Metrológica.} Diseño experimental, sensores NDIR, arquitectura del nodo IoT, protocolo de ventanas cerradas, calibración frente a patrones certificados, compensación cruzada por temperatura y humedad, presupuesto de incertidumbre GUM y resultados de caracterización instrumental.
    
    \item[Capítulo 3] \textbf{Modelo Matemático para la Predicción de Emisiones.} Fundamentos del modelo DEB para \textit{H. illucens}, variables de estado, formulación del sistema de EDO, cinética de asimilación (Holling II), modulación térmica (Sharpe--Schoolfield), parámetros semilla y restricciones físicas, análisis estructural y consistencia del balance de carbono.
    
    \item[Capítulo 4] \textbf{Sistema Híbrido: Integración Modelo--Datos.} Arquitectura del sistema híbrido DEB--MLP, extractor de pendientes en ventanas cerradas, calibrador neuronal parámetrico, estrategia de calibración por fases (semillas Eriksen, ajuste local NDIR, corrección neuronal), evaluación comparativa RNA-B vs. MLP, validación predictiva y resultados de integración modelo--datos.
    
    \item[Capítulo 5] \textbf{Validación Predictiva y Diagnóstico Operativo.} Inventario dinámico de carbono, clasificador de umbral para detección de anaerobiosis, métricas MRV comparables con la literatura, plataforma técnica de soporte a la decisión y semáforo operativo.
    
    \item[Capítulo 6] \textbf{Conclusiones Generales y Perspectivas.} Síntesis de aportes, limitaciones del estudio, trabajos futuros y recomendaciones para la escalabilidad del sistema.
\end{description}

Los anexos técnicos complementan el cuerpo del documento con especificaciones de hardware, formulaciones matemáticas detalladas, algoritmos de aprendizaje automático, datos experimentales completos y comparaciones teóricas con tecnologías alternativas.
